\section*{Hoofdstuk \ref{ch:b1:induction} --- Elektromagnetische inductie en toepassingen}

\begin{solution}{exo:b1:induction:1}
$\Phi = NBS\cos30^\circ = 100 \times 0.2 \times 10^{-3} \times 0.866 = \qty{1.7e-2}{Wb}$;
$\langle e\rangle = \Delta\Phi/\Delta t = \qty{1.7}{V}$ in \qty{10}{ms}, \qty{17}{V} in
\qty{1}{ms}. De stroom loopt zo dat hij de verdwijnende flux herstelt: zijn
veld ligt langs de oude $\vect B$.
\end{solution}

\begin{solution}{exo:b1:induction:2}
$e = B\ell v = 0.5 \times 0.2 \times 2 = \qty{0.20}{V}$; $i = \qty{2.0}{A}$; $F = B\ell i =
\qty{0.20}{N}$ remmend, dus levert wie duwt \qty{0.20}{N}; $Fv = \qty{0.40}{W}
= Ri^2$.
\end{solution}

\begin{solution}{exo:b1:induction:3}
$L = \mu_0n^2S\ell = 4\pi \times 10^{-7} \times 10^6 \times 4 \times 10^{-4} \times 0.2 =
\qty{0.10}{mH}$; $W = \tfrac12LI^2 = \qty{1.3}{mJ}$; $\tau = L/R = \qty{0.1}{ms}$.
$B = \mu_0nI = \qty{6.3}{mT}$, $B^2/2\mu_0 = \qty{15.7}{J/m^3}$ over
$S\ell = \qty{8e-5}{m^3}$: \qty{1.3}{mJ}.
\end{solution}

\begin{solution}{exo:b1:induction:4}
$N_2/N_1 = 12/230 = 0.052$; $N_1 = 50/0.052 = 960$ windingen; $i_1 = 2 \times 0.052 =
\qty{0.10}{A}$. Een gelijkstroom geeft een constante flux: geen $\dd\varphi/\dd t$, geen
secundaire \omterm{def:b1:dc-circuits:voltage}{spanning} --- en de primaire, een kortsluiting, zou verbranden.
\end{solution}

\begin{solution}{exo:b1:induction:5}
(a) De flux (langs de nadering) groeit: de stroom maakt een veld dat
haar tegenwerkt --- zijn eigen noordpool staat tegenover de noordpool van de magneet --- en de
ring wordt weggeduwd. (b) Een stroom waarvan het veld de groei tegenwerkt; de
ring, een dipool die tegengesteld aan $\vect B$ staat, wordt naar het zwakkere
veld geduwd. (c) Bij het vallen groeit de flux: dezelfde stroom als bij (a), de ring
wordt afgeremd. (d) Elk stuk buis vóór de magneet ziet groeiende flux,
erachter afnemende; beide geïnduceerde stromen remmen de magneet; bij de
snelheid waarbij de \omterm{thm:b1:magnetostatics:laplace}{laplacekracht} het gewicht in evenwicht houdt valt de magneet
gestaag en wordt zijn verloren \omterm{def:b1:work-and-energy:conservative}{potentiële energie} joulse warmte in de buis.
\end{solution}

\begin{solution}{exo:b1:induction:6}
$E_{\max} = NBS\omega = 200 \times 2 \times 10^{-3} \times 0.1 \times 314 = \qty{12.6}{V}$, effectief
\qty{8.9}{V}; $I_{\max} = \qty{1.26}{A}$, $\Gamma_{\max} = NBSI_{\max} = \qty{0.050}{N.m}$
(wanneer $\sin\omega t = 1$); $\langle P\rangle = E_{\max}^2/2R = \qty{7.9}{W}$ $=
\langle\Gamma\omega\rangle = \tfrac12 \times 0.05 \times 314$; wanneer $e = 0$ is de stroom
nul en verdwijnt het \omterm{def:b1:angular-momentum:moment}{koppel} --- het \omterm{def:b1:angular-momentum:moment}{koppel} pulseert met \qty{100}{Hz}.
\end{solution}

\begin{solution}{exo:b1:induction:7}
$e = Bav$, $i = Bav/R$, $F = B^2a^2v/R = 10\,v$ (N). $m\,\dd v/\dd t = -B^2a^2v/R$:
$v = v_0\eu^{-t/\tau}$, $\tau = mR/B^2a^2 = \qty{0.10}{s}$; afstand $v_0\tau =
\qty{10}{cm}$ (als het veldgebied lang genoeg is). De \omterm{thm:b1:work-and-energy:ke}{kinetische energie},
\qty{0.5}{J}, wordt als $Ri^2$ in de winding gedissipeerd.
\end{solution}

\begin{solution}{exo:b1:induction:8}
Flux van het veld $\mu_0n_1i_1$ van de solenoïde door de $N_2$ windingen: $M =
\mu_0n_1N_2S_2 = 4\pi \times 10^{-7} \times 2000 \times 50 \times 2 \times 10^{-4} = \qty{25}{\micro H}$.
Oplopend: $e = M\,\dd i_1/\dd t = \qty{2.5}{mV}$; bij \qty{10}{kHz} en \qty{1}{A}: $e_{\max} =
M\omega I = 25 \times 10^{-6} \times 6.3 \times 10^4 = \qty{1.6}{V}$. $M$ is symmetrisch:
hetzelfde getal geeft de flux die de kleine spoel door de
solenoïde stuurt --- veel lastiger rechtstreeks te berekenen.
\end{solution}

\begin{solution}{exo:b1:induction:9}
$W = B^2V/2\mu_0 = 2.25/(2.51 \times 10^{-6}) = \qty{0.90}{MJ}$; $L = 2W/I^2 =
\qty{7.2}{H}$; $h = W/mg = \qty{92}{m}$; helium: $9 \times 10^5/2600 = \qty{350}{L}$.
Aardveld: $(5 \times 10^{-5})^2/2\mu_0 = \qty{1e-3}{J/m^3}$, \qty{50}{mJ} in
de kamer.
\end{solution}

\begin{solution}{exo:b1:induction:10}
(a) $E - B\ell v = Ri$ (de tegen-emk werkt de bron tegen), $m\,\dd v/\dd t =
B\ell i$: $m\,\dd v/\dd t = B\ell(E - B\ell v)/R$. (b) $v_\infty = E/B\ell$, $\tau =
mR/B^2\ell^2$. (c) Vermenigvuldig de elektrische vergelijking met $i$: $Ei = Ri^2 +
B\ell vi$, en $B\ell vi = Fv$ is het mechanische vermogen; \omterm{def:b1:heat-engines:efficiency}{rendement} $Fv/Ei =
B\ell v/E = v/v_\infty$ --- $100\%$ alleen bij nullast, waar niets wordt
geleverd. (d) $v_\infty = 2/(0.5 \times 0.2) = \qty{20}{m/s}$, $\tau = 0.05 \times
0.1/0.01 = \qty{0.5}{s}$. (e) $F = B\ell I = 1 \times 0.1 \times 10^6 = \qty{1e5}{N}$
(met $\ell = \qty{10}{cm}$), $a = \qty{1e7}{m/s^2}$, $v = \sqrt{2a \times 5} =
\qty{1e4}{m/s}$ --- tien kilometer per seconde, de belofte en het
probleem van railkanonnen (de rails eroderen).
\end{solution}

\begin{solution}{exo:b1:induction:11}
(a) Op straal $r$ beweegt het metaal met $v = \omega r$; $\vect v\wedge\vect B$ is
radiaal, met grootte $\omega rB$: $e = \int_0^R\omega Br\,\dd r = \tfrac12B\omega R^2$.
(b) $\omega = 314$: $e = 0.5 \times 1 \times 314 \times 0.01 = \qty{1.6}{V}$. (c) Slechts één
``winding'', dus hooguit volts; in \qty{1}{m\ohm}: \qty{1.6}{kA}, \omterm{def:b1:angular-momentum:moment}{koppel}
$= P/\omega = ei/\omega = 2500/314 = \qty{8}{N.m}$. (d) De kring borstel--straal--
rand--borstel bestrijkt het oppervlak $\tfrac12R^2\omega\,\dd t$ per $\dd t$: $\dd\Phi/\dd t =
\tfrac12BR^2\omega$ --- de flux door de bewegende, vervormende kring.
\end{solution}

\begin{solution}{exo:b1:induction:12}
(a) $e_2 = -M\,\dd i_1/\dd t$, piek $M\omega I_1 = 10^{-6} \times 1.57 \times 10^5 \times 30 =
\qty{4.7}{V}$. (b) $I_2 = 4.7/0.01 = \qty{470}{A}$ piek; $\langle P\rangle = R_pI_2^2/2
= \qty{1.1}{kW}$. (c) De emk is $\propto\omega$: een hoge frequentie geeft een bruikbare
emk bij een bescheiden $M$; het glas is een isolator, er loopt niets in
--- de warmte ontstaat in het metaal zelf. (d) De stroom in de pan induceert in
de spoel $M\omega I_2 = 10^{-6} \times 1.57 \times 10^5 \times 470 = \qty{74}{V}$ piek,
in fase met $i_1$ (de panstroom loopt niet achter op $e_2$, want hij is
resistief, en $e_2$ loopt \qty{90}{\degree}\ldots\ achter op $i_1$, zodat de
teruggekoppelde emk $-M\,\dd i_2/\dd t$ is, in tegenfase met de
\emph{aandrijvende} \omterm{def:b1:dc-circuits:voltage}{spanning} van $i_1$): de generator levert $\tfrac12 \times 74 \times 30 =
\qty{1.1}{kW}$ meer --- precies wat de pan dissipeert.
\end{solution}

\begin{solution}{pb:b1:induction:1}
\textbf{1.} $\Phi = NBS\cos\omega t$, $e = NBS\omega\sin\omega t$; $NBS = 300 \times
0.3 \times 3 \times 10^{-4} = \qty{0.027}{Wb}$; $\omega = v/r = 5/0.01 = \qty{500}{rad/s}$:
\omterm{def:b1:wave-propagation:signal}{amplitude} \qty{13.5}{V}, frequentie $\omega/2\pi = \qty{80}{Hz}$.

\textbf{2.} $I_{\text{rms}} = 13.5/\sqrt2/16 = \qty{0.60}{A}$; lamp $RI^2 = \qty{4.3}{W}$
bij \qty{7.2}{V}: een lamp van \qty{6}{V}, \qty{3}{W} wordt overbelast en gaat niet
lang mee.

\textbf{3.} Mechanisch $=$ elektrisch: $(R + r_c)I_{\text{rms}}^2 = 16 \times 0.36 =
\qty{5.7}{W}$; $F = P/v = 5.7/5 = \qty{1.1}{N}$ --- de helft van de rolweerstand;
de fietser voelt het.

\textbf{4.} Het laplacekoppel van de geïnduceerde stroom werkt de draaiing tegen
(Lenz); de energie van de lamp komt uit de benen van de fietser, \qty{5.7}{W} ervan.

\textbf{5.} $\underline Z = R + r_c + jL\omega$; $I = NBS\omega/\sqrt{(R + r_c)^2 + L^2\omega^2}$;
voor $L\omega \gg R + r_c$ gaat $I \to NBS/L = 0.027/0.02 = \qty{1.35}{A}$. $70\%$: $L\omega
= 0.7(R + r_c)/\sqrt{1 - 0.49} = 15.7$, $\omega = \qty{784}{rad/s}$, $v = \qty{7.8}{m/s}
= \qty{28}{km/h}$.

\textbf{6.} Zonder $L$: $e_{\max} = 3 \times 13.5 = \qty{40.5}{V}$, $I_{\text{rms}} =
\qty{1.8}{A}$, \qty{39}{W} in de lamp --- die knapt. Met $L$: $I = 40.5/\sqrt{256 +
900} = \qty{1.2}{A}$ piek, \qty{8.5}{W}: heet, maar tien keer minder.

\textbf{7.} \qty{9}{km/h}: $e_{\max} = 6.75$, $|Z| = \sqrt{256 + 25} = 16.8$, $I = 0.40$,
$P = \tfrac12 \times 0.16 \times 12 = \qty{1.0}{W}$: zwak. \qty{36}{km/h}: $e_{\max} = 27$,
$|Z| = 25.6$, $I = 1.05$, $P = \qty{6.6}{W}$. Vier keer de snelheid, $2.6$ keer
de stroom: de \omterm{def:b1:induction:inductance}{zelfinductie} regelt.

\textbf{8.} Alleen de onderlinge beweging telt: de flux door de vaste
spoel verandert net zo, dezelfde emk. Er zijn geen sleepcontacten (borstels)
nodig om de stroom uit een draaiende spoel te halen.

\textbf{9.} Geen beweging, geen fluxverandering, geen licht --- vandaar de \omterm{def:b1:potential-capacitors:capacitance}{condensator} of
kleine accu (``standlicht'') die tijdens het rijden wordt geladen.

\textbf{10.} Elk draadelement staat loodrecht op het radiale $\vect B$:
$F = B\ell i$, axiaal. Bewegend met $v$ ziet elk element $\vect v\wedge\vect B$
langs de draad: $e = -B\ell v$, tegengesteld aan de stroom die de
beweging veroorzaakt (verbruikersconventie: een \omterm{def:b1:dc-circuits:voltage}{spanning} $B\ell v$).

\textbf{11.} $u = Ri + L\,\dd i/\dd t + B\ell v$; $m\,\dd v/\dd t = -kx - hv + B\ell i$.

\textbf{12.} $\underline i = (\underline u - B\ell\underline v)/(R + jL\omega)$; $\underline v\,(jm\omega + h +
k/j\omega) = B\ell\,\underline i$; vul in en los op naar $\underline v$.

\textbf{13.} $i \approx (u - B\ell v)/R$ geeft $m\dot v + (h + B^2\ell^2/R)v + kx =
B\ell u/R$: extra demping $25/8 = \qty{3.1}{N.s/m}$. $f_0 = \tfrac1{2\pi}\sqrt{k/m} =
\qty{50}{Hz}$; $Q = \sqrt{km}/h_{\text{tot}}$: $3.2$ open, $0.77$ met de
versterker.

\textbf{14.} $F = \qty{5.0}{N}$, $a = F/m = \qty{500}{m/s^2}$, $x = a/\omega^2 =
\qty{13}{\micro m}$, $v = a/\omega = \qty{0.080}{m/s}$; $B\ell v = \qty{0.40}{V}$ tegen
$Ri = \qty{8}{V}$: $5\%$.

\textbf{15.} Met opgelegde stroom is $v = F/h = \qty{5}{m/s}$ (een
\omterm{def:b1:wave-propagation:signal}{amplitude} van \qty{16}{mm} --- een grote slag voor \qty{1}{A}); $B\ell v =
\qty{25}{V}$, drie keer $Ri$: de versterker moet \qty{33}{V} leveren om
\qty{1}{A} bij resonantie door te drukken --- de impedantiepiek. Een versterker als \omterm{def:b1:dc-circuits:sources}{spanningsbron}
zou de elektrische demping de conus laten vasthouden.

\textbf{16.} $\tfrac12RI^2 = \qty{4.0}{W}$; $\tfrac12hv^2 = \tfrac12 \times 0.0064 =
\qty{3.2}{mW}$: \omterm{def:b1:heat-engines:efficiency}{rendement} $0.08\%$ --- luidsprekers zijn kachels die
fluisteren.

\textbf{17.} Boven de resonantie is $a = F/m = B\ell i/m$ bij gegeven stroom
onafhankelijk van de frequentie, zodat de druk, $\propto a$, vlak is:
de massagestuurde band is de werkband.

\textbf{18.} Radiaal veld: de kracht is axiaal en overal in de spleet even
groot, waar de spoel ook staat. Een tweeter moet tot \qty{20}{kHz}
massagestuurd blijven met een minuscule slag: lichte spoel en conus; en een
kleine conus straalt hoge tonen naar alle kanten uit.

\textbf{19.} Lage frequentie: $v$ klein, $|Z| \approx R = \qty{8}{\ohm}$; bij
\qty{50}{Hz}: $u = Ri + B\ell v$ met $v = B\ell i/h$: $Z = R + B^2\ell^2/h = 8 + 25 =
\qty{33}{\ohm}$; bij \qty{10}{kHz}: $|R + jL\omega| = |8 + 31j| = \qty{32}{\ohm}$. Een
piek bij resonantie, een vlakke bodem van \qty{8}{\ohm}, en een stijging met $L\omega$.

\textbf{20.} $e = B\ell v = 5 \times 10^{-3} = \qty{5}{mV}$.

\textbf{21.} $i = 5 \times 10^{-3}/608 = \qty{8.2}{\micro A}$; $F = B\ell i = \qty{41}{\micro N}$,
tegengesteld aan $v$: $F = [B^2\ell^2/(R + R_L)]\,v$, een demping van \qty{0.041}{N.s/m};
het vermogen $Fv = \qty{4e-8}{W}$ komt uit het geluid en belandt in de
twee weerstanden.

\textbf{22.} Spoeldrager: $B^2\ell_f^2/R_f = 1 \times 0.01/10^{-3} = \qty{10}{N.s/m}$, tien
keer $h$: hij zou de conus verdoven en vermogen als warmte verspillen; een spleet
onderbreekt de gesloten winding.

\textbf{23.} Open: alleen $h$ dempt, $Q = 3.2$, hij natrilt. Kortgesloten: de
tegen-emk drijft een stroom door $R$, de demping $B^2\ell^2/R$ komt erbij,
$Q = 0.77$: hij stopt meteen. (Probeer het.)

\textbf{24.} Over een periode keren de opgeslagen energieën naar hun waarde terug:
$\langle ui\rangle = \langle Ri^2\rangle + \langle hv^2\rangle$ --- elektrische toevoer $=$
joulse warmte in de spoel $+$ mechanische verliezen (uitgestraald geluid en
wrijving in de ophanging); de koppelterm $B\ell iv$ komt in beide
vergelijkingen met tegengesteld teken voor en valt weg.

\textbf{25.} Dynamo: Faraday op een draaiende spoel --- \qty{13.5}{V} bij
\qty{18}{km/h}, met de stroom door haar eigen \omterm{def:b1:induction:inductance}{zelfinductie} afgetopt op \qty{1.35}{A}.
Luidspreker: \omterm{thm:b1:magnetostatics:laplace}{laplacekracht} en tegen-emk --- \qty{3.1}{N.s/m} elektrische
demping, $0.08\%$ \omterm{def:b1:heat-engines:efficiency}{rendement}. Microfoon: dezelfde spoel achterstevoren ---
\qty{5}{mV} per millimeter per seconde.
\end{solution}
